\documentclass[11pt]{article}
\usepackage{docstyle}
\renewcommand\sheetname{Light \textperiodcentered{} Homework answers}
\FigAnstrue

\begin{document}
\sheethead{Light: Unit Review Homework \textperiodcentered{} Answers}{Tutor's working, not an official mark scheme}{Answers only}
Drawings show the given parts in black and the answer in \ans{red}. For Section B, each answer says what reaches Achieved (A), Merit (M) and Excellence (E); the full answer given is the Excellence answer.

\section*{Section A}
\A{A1} the Sun \textbf{S} \quad the Moon \textbf{R} \quad a candle flame \textbf{S} \quad a bicycle reflector \textbf{R} \quad a mirror \textbf{R} \quad a glow-worm \textbf{S}

\A{A2} frosted glass \textbf{translucent} \quad clear plastic bottle \textbf{transparent} \quad aluminium foil \textbf{opaque} \quad tracing paper \textbf{translucent}

\A{A3} time $= \dfrac{\text{distance}}{\text{speed}} = \dfrac{384\,000\ \mathrm{km}}{300\,000\ \mathrm{km/s}} = \mathbf{1.28\ s}$\par
Check: $1.28 \times 300\,000 = 384\,000$. An answer of about 0.78 means the division was the wrong way round.

\begin{minipage}[t]{0.5\textwidth}
\vspace{0pt}
\A{A4} (a) Normal at 90\textdegree{} to the mirror at the point where the ray hits; reflected ray leaves at 20\textdegree{} to the mirror on the other side of the normal.\par
(b) angle of incidence $= 90^\circ - 20^\circ = \mathbf{70^\circ}$; angle of reflection $= \mathbf{70^\circ}$.\par
Answering 20\textdegree{} measures from the mirror (L1).
\end{minipage}\hfill
\begin{minipage}[t]{0.46\textwidth}
\vspace{0pt}\centering\TaskReflect[5.8mm]{1}
\end{minipage}

\A{A5} (a) \textbf{2.5\,m} behind the mirror.\par
(b) The student is now $2.5 - 1.0 = 1.5$\,m from the mirror, and the image is 1.5\,m behind it: $1.5 + 1.5 = \mathbf{3.0\,m}$.

\A{A6} green leaf in white light: \textbf{green} \quad white T-shirt in blue light: \textbf{blue} \quad red rose in green light: \textbf{black}

\A{A7} (a) \textbf{towards the normal} (light slows down) \quad (b) \textbf{away from the normal} (light speeds up) \quad (c) \textbf{no bending} (along the normal)

\A{A8} (a) \textbf{iris} \quad (b) \textbf{retina} \quad (c) \textbf{optic nerve} \quad (d) \textbf{cornea}

\section*{Section B}
\A{B1} The flower looks \textbf{black}.\par
A yellow object \textbf{reflects} yellow light, and yellow light is a mix of \textbf{red and green}, so the flower reflects red and green and \textbf{absorbs} all other colours. Blue light contains \textbf{no red or green} light. So the flower absorbs all the blue light and \textbf{no light is reflected into our eyes}, and the flower appears black.\par
\emph{A}: states black, or says the blue light is absorbed. \emph{M}: says yellow objects reflect yellow (red and green) and absorb other colours. \emph{E}: the complete chain, ending with no light reaching the eye.

\begin{minipage}[t]{0.52\textwidth}
\vspace{0pt}
\A{B2} (a) Four straight rays with arrows from the top and bottom of the lamp past the top and bottom of the puppet. The dark centre is the \textbf{umbra}; the lighter regions above and below are the \textbf{penumbra}.
\end{minipage}\hfill
\begin{minipage}[t]{0.44\textwidth}
\vspace{0pt}\centering\TaskPuppet[5.2mm]
\end{minipage}\par
(b) The lamp is \textbf{large}. Light travels in straight lines, so at the edge of the shadow the puppet blocks light from \textbf{only part} of the lamp, while light from the rest of the lamp still reaches the screen. This partly lit region is the \textbf{penumbra}. Moving outwards across it, more and more of the lamp is visible, so the shadow fades gradually: the edge looks fuzzy.\par
(c) The shadow gets \textbf{bigger}. Light travels in straight lines from the lamp past the edges of the puppet. When the puppet is closer to the lamp, these rays leave the puppet at a \textbf{wider angle}, so they are further apart when they reach the screen.\par
\emph{A}: correct drawing, or ``bigger''. \emph{M}: reason with straight lines, or ``only some of the light is blocked''. \emph{E}: both explanations linked to rays from a large source travelling in straight lines.

\begin{minipage}[t]{0.52\textwidth}
\vspace{0pt}
\A{B3} (a) Ray from the coin to the surface, bending \textbf{away from the normal} into the eye, with arrows towards the eye; the ray in air extended back as a dotted line to the \textbf{image}, which is above the coin.\par
\end{minipage}\hfill
\begin{minipage}[t]{0.44\textwidth}
\vspace{0pt}\centering\TaskFish[5.4mm]{Po}
\end{minipage}\par
(b) Light from the coin travels from water into air. It \textbf{speeds up}, so it is refracted \textbf{away from the normal} at the surface. Our brain assumes light travels in \textbf{straight lines}, so it traces the ray straight back. The ray appears to come from a point \textbf{higher} than the coin, so we see a \textbf{virtual image} of the coin closer to the surface, and the pool looks shallower than it is.\par
\emph{A}: ``the light is refracted''. \emph{M}: speeds up, bends away from the normal. \emph{E}: also the straight-line tracing back to a higher virtual image.

\A{B4} \textbf{Normal eye:} light from the distant object enters through the \textbf{cornea} and the \textbf{pupil}. The cornea and the convex \textbf{lens} \textbf{refract} the light so that the rays \textbf{converge}. They meet exactly on the \textbf{retina}, forming a clear, real image. The retina sends signals along the \textbf{optic nerve} to the brain.\par
\textbf{Short sight:} the eye converges the light \textbf{too strongly} (the lens is too strong or the eyeball too long), so rays from a distant object meet \textbf{in front of} the retina. By the time they reach the retina they have spread out again, so the image is blurred.\par
\textbf{Concave lens:} a concave lens makes the rays \textbf{diverge} (spread out) slightly before they enter the eye. The eye then brings them together \textbf{further back}, so the focus moves onto the retina and the image is clear.\par
\emph{A}: names the parts, or ``focus in front of the retina''. \emph{M}: explains converging on the retina, or why a concave lens is used. \emph{E}: all three parts linked: converge on retina; converge too much so focus in front; concave lens diverges so the focus moves back onto the retina.

\A{B5} A magenta filter \textbf{transmits red and blue} light and \textbf{absorbs green}. So the dancer is lit with red and blue light only.
\begin{itemize}
\item White shirt: reflects all colours, so it reflects the red and blue it receives and looks \textbf{magenta}.
\item Green trousers: reflect only green; there is \textbf{no green} in the light, so the red and blue are absorbed and they look \textbf{black}.
\item Blue hat: reflects blue and absorbs red, so it looks \textbf{blue}.
\end{itemize}
\emph{A}: one or two correct colours. \emph{M}: all three colours with the filter explained (transmits red and blue). \emph{E}: each colour explained by what the object reflects and absorbs from the light that reaches it.

\A{B6} White light is a \textbf{mixture} of colours. When it enters the glass, the light \textbf{slows down} and is refracted. Each colour travels at a \textbf{different speed} in glass, so each colour is refracted by a \textbf{different amount}. Violet is slowed the most, so it is refracted the \textbf{most}; red is slowed the least, so it is refracted the \textbf{least}. The colours separate, and they separate further when they are refracted again at the second face, so violet and red end up on opposite sides of the spectrum. This splitting is called \textbf{dispersion}.\par
\emph{A}: ``the prism splits (disperses) white light''. \emph{M}: different colours are refracted by different amounts. \emph{E}: linked to different speeds, with violet bent most and red least.

\end{document}
